% **************************************************************************** %
\chapter{\code{LTspice}-Schaltungen}
\label{app:ltspice}
% **************************************************************************** %

Dieses Kapitel  beinhaltet \code{LTspice}-Schaltungen, welche  zu Simulationen
benutzt worden sind. Erkl\"arungen  zu den jeweiligen Schaltungen  sind in den
Kapiteln  zu finden,  welche auf  sie verweisen. S\"amtliche  Schaltungen sind
auch elektronisch als \code{.asc}-Datei verf\"ugbar.

\noindent\begin{minipage}{\textwidth}
    \centering
    \includegraphics[width=0.9\textwidth]{images/ltspice/singlecell.eps}
    \figcaption{%
        Schaltung        zur         Simulation        einer        Solarzelle
        gem\"ass        Abbildung       \ref{fig:circuit:solarCell}        von
        Seite     \pageref{fig:circuit:solarCell}. Die     Solarmodule     aus
        den     Abbildungen    \ref{fig:ltspice:module:cellBased:36x2}     und
        \ref{fig:ltspice:module:cellBased:72x1}  basieren auf  Zellen gem\"ass
        diesem Schaltschema.%
    }
    \label{fig:ltspice:solarCell}
\end{minipage}

\begin{figure}[h!tb]
    \centering
    \includegraphics[width=\textwidth]{images/ltspice/module-72cells--72x1.eps}
    \caption{%
        Solarmodul   \"ahnlich   zu   Sunmodule  Pro-Series   XL   SW320   aus
        \cite{ref:solar:sunmodulePro},  modelliert durch  einen Strang  mit 72
        Zellen gem\"ass Abbildung \ref{fig:ltspice:solarCell} in Serie.%
    }
    \label{fig:ltspice:module:cellBased:72x1}
\end{figure}

\begin{figure}[h!tb]
    \centering
    \includegraphics[angle=90,width=\textwidth]{images/ltspice/module-string-complete.eps}
    \caption{%
        Schaltung  zur  Erzeugung  von  Strom-Spannungs-Kurven  mit  separatem
        Leitungs-  und  Lastwiderstand. Diese  Schaltung ist  benutzt  worden,
        um  die  Kurven  in   Abbildung  \ref{fig:ivcurve:modules}  auf  Seite
        \pageref{fig:ivcurve:modules} zu erzeugen.%
    }
    \label{fig:ltspice:vicurvesimulator:modules}
\end{figure}
